%% Font size %%
\documentclass[11pt]{article}

%% Load the custom package
\usepackage{Mathdoc}

%% Numéro de séquence %% Titre de la séquence %%
\renewcommand{\centerhead}{S14- Probabilités - Exercices Divers}

%% Spacing commands %%
\renewcommand{\baselinestretch}{1} \setlength{\parindent}{0pt}

%%\pagestyle{empty}

\begin{document}

\begin{exercice}[3][Problème — Tombola]

Lors d'une tombola, 200 tickets sont vendus.

Parmi eux :

\begin{itemize}
    \item 1 ticket permet de gagner un vélo ;
    \item 4 tickets permettent de gagner un ballon ;
    \item 20 tickets permettent de gagner une peluche.
\end{itemize}

On choisit un ticket au hasard.

\begin{enumerate}
    \item Quelle est la probabilité de gagner le vélo ?
    \item Quelle est la probabilité de gagner un ballon ?
    \item Quelle est la probabilité de gagner un cadeau ?
    \item Quelle est la probabilité de ne rien gagner ?
\end{enumerate}

\end{exercice}


\begin{exercice}[3][Problème — Météo]

Un professeur note la météo pendant 30 jours.

\begin{center}
\begin{tabular}{|c|c|}
\hline
Temps & Nombre de jours \\
\hline
Soleil & 12 \\
Nuages & 10 \\
Pluie & 8 \\
\hline
\end{tabular}
\end{center}

On choisit un jour au hasard.

\begin{enumerate}
    \item Calculer la probabilité qu'il ait fait soleil.
    \item Calculer la probabilité qu'il ait plu.
    \item Calculer la probabilité qu'il n'ait pas plu.
\end{enumerate}

\end{exercice}


\newpage


\begin{exercice}[3][Tableau de probabilités]

Une urne contient des jetons rouges, bleus et verts.

Compléter le tableau.

\begin{center}
\begin{tabular}{|c|c|c|}
\hline
Couleur & Nombre de jetons & Probabilité \\
\hline
Rouge & 4 & \\
\hline
Bleu & 3 & \\
\hline
Vert & 5 & \\
\hline
Total & & 1 \\
\hline
\end{tabular}
\end{center}

\end{exercice}


\begin{exercice}[3][Comparer des probabilités]

Dans un premier sac, il y a 2 boules rouges et 3 boules bleues.

Dans un deuxième sac, il y a 5 boules rouges et 10 boules bleues.

\begin{enumerate}
    \item Calculer la probabilité d'obtenir une boule rouge dans chaque sac.
    \item Dans quel sac a-t-on le plus de chances d'obtenir une boule rouge ?
\end{enumerate}

\end{exercice}


\begin{exercice}[3][Chercher une quantité]

Dans un sac contenant 20 boules, la probabilité d'obtenir une boule rouge est $\dfrac{7}{20}$.

\begin{enumerate}
    \item Combien y a-t-il de boules rouges ?
    \item Combien y a-t-il de boules qui ne sont pas rouges ?
\end{enumerate}

\end{exercice}

\begin{exercice}[3][Jeu télévisé]

Dans un jeu, un candidat choisit au hasard une enveloppe parmi 50.

\begin{itemize}
    \item 5 enveloppes contiennent 100 € ;
    \item 15 enveloppes contiennent 20 € ;
    \item les autres enveloppes sont vides.
\end{itemize}

\begin{enumerate}
    \item Quelle est la probabilité de gagner 100 € ?
    \item Quelle est la probabilité de gagner de l'argent ?
    \item Quelle est la probabilité de ne rien gagner ?
\end{enumerate}

\end{exercice}


\begin{exercice}[4][Exercice de recherche]

Un jeu consiste à lancer un dé à 6 faces.

\begin{itemize}
    \item Si le joueur obtient 1 ou 2, il gagne 2 points.
    \item S'il obtient 3, 4 ou 5, il gagne 1 point.
    \item S'il obtient 6, il perd 1 point.
\end{itemize}

\begin{enumerate}
    \item Quelle est la probabilité de gagner 2 points ?
    \item Quelle est la probabilité de gagner au moins 1 point ?
    \item Quelle est la probabilité de perdre 1 point ?
\end{enumerate}

\end{exercice}


\begin{exercice}[4][Créer une expérience aléatoire]

Inventer une expérience aléatoire.

Puis :

\begin{enumerate}
    \item donner toutes les issues possibles ;
    \item proposer deux événements ;
    \item calculer la probabilité d'un des événements.
\end{enumerate}

\end{exercice}

\end{document}
